
\documentclass[11pt]{article}
\usepackage[margin=1in]{geometry}
\usepackage{amsmath,amssymb,booktabs,longtable,array,enumitem,hyperref,siunitx,microtype}
\usepackage[numbers,sort&compress]{natbib}
\hypersetup{colorlinks=true,linkcolor=blue,citecolor=blue,urlcolor=blue}

\title{Technical Formulation and Parameter Justification for the Variable-Compliance Rover Mobility Model}
\author{Jane Park}
\date{\today}

\begin{document}
\maketitle

\section{Purpose}

This document records the mathematical formulation used in the variable-compliance rover mobility model and identifies the justification and provenance of each term. It is intended to accompany the research report and source code.

\section{State variable}

The wheel state is represented by
\begin{equation}
s\in[0,1].
\end{equation}

Recovered source anchor states are
\begin{equation}
s_k=
\{0,\;0.149537,\;0.422427,\;0.572765,\;0.746899,\;0.916081,\;1\}.
\end{equation}

The continuous constitutive response is
\begin{equation}
D=D(F,s),
\end{equation}
constructed from Lee et al.\ force--deformation data \cite{Lee2024}.

\paragraph{Justification.}
A continuous state avoids imposing an artificial binary soft/stiff wheel and permits optimization over the experimentally represented compliance range.

\section{Normal load}

For the present single-wheel-equivalent calibration,
\begin{equation}
F_N=W=mg.
\end{equation}

Using
\begin{equation}
m=5.3~\mathrm{kg},\qquad g=9.81~\mathrm{m\,s^{-2}},
\end{equation}
gives
\begin{equation}
F_N=51.993~\mathrm{N}.
\end{equation}

For a general $N_w$-wheel rover under a static equal-load approximation,
\begin{equation}
F_{N,k}\approx\frac{m_{\mathrm{rover}}g}{N_w}.
\end{equation}

\paragraph{Justification.}
The current fixed value preserves consistency with the Lee small-wheel data. A general rover implementation should use per-wheel load rather than total rover weight.

\section{Wheel geometry}

The current calibration uses
\begin{equation}
R=0.140~\mathrm{m},
\qquad
b_w=0.040~\mathrm{m}.
\end{equation}

\paragraph{Justification.}
These dimensions are tied to the source wheel configuration used to define $D(F,s)$. They are not yet free wheel-family parameters.

\section{Contact geometry}

For radial deformation $D$,
\begin{equation}
L_c=2\sqrt{2RD-D^2}.
\end{equation}

The reduced contact area is
\begin{equation}
A=b_wL_c,
\end{equation}
and the effective Bekker width is
\begin{equation}
b=\min(b_w,L_c).
\end{equation}

Average pressure is
\begin{equation}
p=\frac{F_N}{A}.
\end{equation}

\paragraph{Justification.}
This reduced geometry allows the Lee-derived deformation response to modify the contact patch used by the terramechanics calculation. It is an approximation and does not represent the full pressure distribution.

\section{Pressure--sinkage relation}

For deformable terrain,
\begin{equation}
\sigma=
\left(
\frac{k_c}{b}+k_\phi
\right)y^n,
\end{equation}
following classical Bekker/Wong terramechanics \cite{Bekker1969,Wong2010}.

Using the average pressure $p$,
\begin{equation}
z=
\left[
\frac{p}
{k_c/b+k_\phi}
\right]^{1/n}.
\end{equation}

\paragraph{Parameters.}
$k_c$, $k_\phi$, and $n$ are terrain-specific pressure--sinkage parameters.

\paragraph{Justification.}
These parameters are terrain properties rather than wheel fitting coefficients. They must be traced to their literature or inherited data source.

\section{Compaction resistance}

The reduced compaction resistance is
\begin{equation}
R_c=
b_w
\left(
\frac{k_c}{b}+k_\phi
\right)
\frac{z^{n+1}}{n+1}.
\end{equation}

\paragraph{Justification.}
This term is a reduced work-per-distance representation of soil compaction. It provides a state-dependent mobility cost because compliance changes contact pressure and sinkage.

\section{Traction feasibility}

Available traction is bounded by
\begin{equation}
F_{\max}=cA+F_N\tan\phi,
\end{equation}
where $c$ is cohesion and $\phi$ is internal friction angle.

The state is rejected if
\begin{equation}
F_{\mathrm{req}}>F_{\max}.
\end{equation}

\paragraph{Justification.}
This is used only as a feasibility bound. It is not interpreted as a calibrated slip prediction. A Janosi--Hanamoto slip-energy model was not added because additional wheel-specific calibration would be required.

\section{Slope}

For a directed graph edge $i\rightarrow j$ with physical length $L_{ij}$,
\begin{equation}
\theta_{ij}
=
\tan^{-1}
\left(
\frac{z_j-z_i}{L_{ij}}
\right).
\end{equation}

The gravitational contribution is
\begin{equation}
F_g=W\sin\theta_{ij}.
\end{equation}

\paragraph{Justification.}
Signed slope is computed directly from elevation and physical edge length. An uphill regression reproduced the expected gravitational energy increment $W\Delta z$.

\section{Positive-propulsion objective}

The raw mobility demand is converted to positive propulsion demand as
\begin{equation}
F_{\mathrm{prop}}=\max(0,F_{\mathrm{raw}}).
\end{equation}

\paragraph{Justification.}
The model does not assign regenerative credit or braking losses. Therefore downhill results must be interpreted as positive propulsion demand only.

\section{Rolling-loss sensitivity term}

The reduced state-dependent rolling term is
\begin{equation}
R_{\mathrm{roll}}(s)
=
W
\left[
C_{rr,\mathrm{ref}}+\lambda h(s)
\right].
\end{equation}

The model retains four alternative $h(s)$ forms:
\begin{enumerate}[label=(\arabic*)]
\item linear state,
\item square-root loading work,
\item loading work,
\item deflection.
\end{enumerate}

The tested sensitivity set uses
\begin{equation}
C_{rr,\mathrm{ref}}
\in\{0.0005,0.001,0.002,0.005\},
\end{equation}
and
\begin{equation}
\lambda\in\{0.003,0.01,0.02\}.
\end{equation}

\paragraph{Justification.}
Lee's quasi-static deformation data do not calibrate cyclic rolling dissipation. The rolling term is therefore treated as model-form and amplitude sensitivity rather than a measured law.

\section{Terrain cost density}

For each terrain $T$, slope $\theta$, and wheel state $s$, the model evaluates
\begin{equation}
c(T,\theta,s)
\end{equation}
in units of positive propulsion work per distance.

For deformable terrain this includes the reduced soil-compaction term, rolling term, signed gravitational contribution, and traction feasibility.

\section{Graph edge length}

For physical grid spacing $\Delta x$,
\begin{equation}
L_{ij}=
\begin{cases}
\Delta x, & \text{cardinal edge},\\
\sqrt{2}\Delta x, & \text{diagonal edge}.
\end{cases}
\end{equation}

\paragraph{Justification.}
$\Delta x$ is the spatial resolution of the map. It is not fitted. A refinement audit demonstrated first-order convergence of adaptive terrain-interface error.

\section{Edge cost}

For an edge connecting terrain endpoints $i$ and $j$,
\begin{equation}
C_{ij}(s)
=
L_{ij}
\frac{
c_i(s,\theta_{ij})+
c_j(s,\theta_{ij})
}{2}.
\end{equation}

\paragraph{Justification.}
Endpoint averaging replaces destination-only cost assignment and reduces directional bookkeeping bias at terrain boundaries.

\section{Grid connectivity}

The final 2-D formulation uses 8-neighbor connectivity.

\paragraph{Justification.}
A 4-neighbor grid produced a resolution-independent 41.42\% path-length error for \SI{45}{\degree} travel. The 8-neighbor formulation removes that error at \SI{45}{\degree}. Residual maximum angular path-length error is approximately 8.24\% near \SI{22.5}{\degree}.

\section{Obstacle feasibility}

For an explicitly defined directed obstacle edge of tested height $h$,
\begin{equation}
s_e\ge s_{\min}(h).
\end{equation}

\paragraph{Justification.}
Only exact Lee-tested heights are supported. The obstacle data define passability, not energy.

\section{Fixed optimization}

For a fixed state $s$ and path $P$,
\begin{equation}
J(P,s)
=
\sum_{e\in P}C_e(s).
\end{equation}

The globally optimal fixed solution is
\begin{equation}
J_{\mathrm{fixed}}^*
=
\min_s
\min_P
J(P,s).
\end{equation}

\paragraph{Justification.}
This establishes a strong baseline. The adaptive wheel is not compared against an arbitrarily selected fixed compliance.

\section{Adaptive optimization}

For an adaptive state sequence $\{s_e\}$,
\begin{equation}
J_{\mathrm{adaptive}}^*
=
\min_{P,\{s_e\}}
\sum_{e\in P}
C_e(s_e).
\end{equation}

Transition energy is currently
\begin{equation}
E_{\mathrm{switch}}=0.
\end{equation}

\paragraph{Justification.}
Actuation energy and transition distance are not sufficiently known. Setting transition energy to zero deliberately defines an ideal mobility benefit rather than inventing an unsupported penalty.

\section{Ideal mobility benefit}

Define
\begin{equation}
\rho
=
\frac{J_{\mathrm{adaptive}}^*}
{J_{\mathrm{fixed}}^*}.
\end{equation}

Then
\begin{equation}
\boxed{
B_{\mathrm{ideal}}
=
1-\rho
=
\frac{
J_{\mathrm{fixed}}^*
-
J_{\mathrm{adaptive}}^*
}{
J_{\mathrm{fixed}}^*
}
}.
\end{equation}

Since the adaptive optimizer can reproduce any fixed state,
\begin{equation}
J_{\mathrm{adaptive}}^*
\le
J_{\mathrm{fixed}}^*,
\end{equation}
and therefore
\begin{equation}
B_{\mathrm{ideal}}\ge0.
\end{equation}

\section{Benefit decomposition}

Let $J_{a|f}$ denote adaptive state selection constrained to the best fixed route. Then
\begin{equation}
\Delta J_{\mathrm{state}}
=
J_f-J_{a|f},
\end{equation}
\begin{equation}
\Delta J_{\mathrm{route}}
=
J_{a|f}-J_a,
\end{equation}
and
\begin{equation}
\Delta J_{\mathrm{total}}
=
\Delta J_{\mathrm{state}}
+
\Delta J_{\mathrm{route}}.
\end{equation}

\paragraph{Justification.}
This separates savings caused by changing wheel state from savings caused by selecting a different route.

\section{Parameter/provenance table}

\begin{longtable}{p{0.17\linewidth}p{0.17\linewidth}p{0.22\linewidth}p{0.36\linewidth}}
\toprule
Quantity & Current value/range & Classification & Justification / status\\
\midrule
$m$ & 5.3 kg & source-matched nominal & Lee small-wheel anchored; not universal rover mass.\\
$F_N$ & 51.993 N & derived & $mg$ for current calibration; general rover should use per-wheel load.\\
$R$ & 0.140 m & source geometry & Not yet a free wheel-family parameter.\\
$b_w$ & 0.040 m & source geometry & Same limitation as $R$.\\
$s$ & 0--1 & source-derived & Recovered continuous compliance coordinate.\\
$D(F,s)$ & recovered surface & source-derived & Main wheel constitutive input.\\
$k_c,k_\phi,n$ & terrain-specific & literature/inherited & Pressure--sinkage properties.\\
$c,\phi$ & terrain-specific & literature/inherited & Traction-feasibility properties.\\
$C_{rr,\mathrm{ref}}$ & swept & sensitivity & Not a measured Lee-wheel value.\\
$\lambda$ & swept & sensitivity & State-dependent rolling-loss amplitude.\\
$h(s)$ & four forms & model-form sensitivity & Explicit uncertainty treatment.\\
$\Delta x$ & map-dependent & physical/numerical & Physical map resolution; convergence audited.\\
Connectivity & 8-neighbor & numerical choice & Selected after quantified Manhattan bias.\\
$E_{\mathrm{switch}}$ & 0 & idealization & Transition cost unknown.\\
Obstacle energy & none & excluded & Data support feasibility only.\\
Downhill regen & none & reduced objective & Negative propulsion clipped to zero.\\
\bottomrule
\end{longtable}

\section{Generality boundary}

The graph and optimization formulation can accept different wheel loads, geometries, terrain maps, and route configurations in principle. The current physical calibration cannot yet support arbitrary wheel geometry.

A broader wheel-family model requires either:
\begin{enumerate}[label=(\arabic*)]
\item experimental data at multiple wheel sizes and loads;
\item validated FEA across geometry/load space; or
\item a mechanics-based dimensionless scaling law.
\end{enumerate}

Until such validation is available, changes in $R$ or $b_w$ should not be interpreted as validated predictions.

\begin{thebibliography}{9}
\bibitem{Lee2024}
J.-Y. Lee et al., ``Variable-stiffness--morphing wheel inspired by the surface tension of a liquid droplet,'' \emph{Science Robotics}, vol. 9, no. 93, eadl2067, 2024. doi: 10.1126/scirobotics.adl2067.

\bibitem{Bekker1969}
M. G. Bekker, \emph{Introduction to Terrain-Vehicle Systems}. University of Michigan Press, 1969.

\bibitem{Wong2001}
J. Y. Wong, \emph{Theory of Ground Vehicles}, 3rd ed. John Wiley \& Sons, 2001.

\bibitem{Wong2010}
J. Y. Wong, \emph{Terramechanics and Off-Road Vehicle Engineering: Terrain Behaviour, Off-Road Vehicle Performance and Design}, 2nd ed. Elsevier, 2010.

\bibitem{Shen2024}
Y. Shen, M. Zou, H. Cao, D. Pan, B. Yuan, and L. He, ``Current-Based Analysis and Validation of a Wheel--Soil Interaction Model for the Trafficability of a Planetary Rover,'' \emph{Aerospace}, vol. 11, no. 11, p. 892, 2024. doi: 10.3390/aerospace11110892.
\end{thebibliography}
\end{document}
